\begin{papercard}{arora2022}
\cf{Stages} \stagetag{C} (GDPR and CCPA/CPRA annotation scheme for app policies), \stagetag{I} (Fleiss' $\kappa$ per category and language among trained annotators), \stagetag{L} (the MAPP corpus later used to evaluate LLMs by Rodriguez et al.\ \pp{rodriguez2024}). LREC 2022 paper (CMU, Bochum, Münster, Penn State, Fordham). Added from the gap search of 23 October 2026 because it holds the only annotated German app-policy corpus found and compares German and English versions of the same apps.

\cf{Corpus and collection} Mobile-app policies from eight Google Play categories. German policies were found through a German IP address and browser locale \pc{arora2022}{3}{To more efficiently find apps with German policies, we relied on the “Recommended for you” Google Play Store feature and configured a server with a German IP address and browser locale}. A policy counted as ``GDPR-aware'' if it used GDPR key phrases or referred to the Privacy Shield; 63 English and 74 German policies of 116 app pairs qualified \pc{arora2022}{3}{Out of our policy pairs for 116 apps, 63 English policies and 74 German policies were found to be “GDPR-aware.”}.

\cf{Annotation} Teams of law students in each language, with the scheme in both languages \pc{arora2022}{4}{To manually annotate the privacy policies, we recruited teams of law students in Germany and the US, who respectively had strong fluency in German and En- glish}; 10 English and 12 German annotators, 1 h 52 min per policy \pc{arora2022}{4}{Annotators spent an average of 1 hour and 52 minutes per policy}. Weekly meetings on hard cases; three annotations per policy; 91 German and 64 English policies completed, 59 in both languages (MAPP-59) \pc{arora2022}{4}{we ultimately obtained 91 German and 64 English fully annotated privacy policies, i.e., with three annotations each, with an intersection of 59 policies}.

\cf{Agreement} Fleiss' $\kappa$ at segment level, per language and level of the scheme \pc{arora2022}{5}{We then computed Fleiss’ Kappa agreement val-}. Category $\kappa$ (Table 4), English / German: first party 0.61 / 0.52, third party 0.52 / 0.47, information type 0.54 / 0.48, purpose 0.63 / 0.58, collection process 0.44 / 0.33, legal basis 0.37 / 0.39, third-party entity 0.49 / 0.36. Classifiers were trained only where $\kappa$ exceeded 0.3 \pc{arora2022}{5}{we focus on building classi- fiers to automatically detect attributes and values for which we had sufficient agreement between anno- tators (Fleiss’ Kappa > 0.3)}.

\cf{German and English versions differ} Two bilingual authors compared the MAPP-59 pairs (German policies machine-translated for the comparison). Fourteen apps acknowledged GDPR in the German policy but not in the English one \pc{arora2022}{7}{A total of 14 apps in MAPP-59 acknowledge GDPR in their German pri- vacy policy but not in the corresponding English pol- icy}; 13 pairs were substantially different and 9 nearly identical \pc{arora2022}{8}{Our manual inspection of the policy pairs in MAPP-59 found that 13 pairs have substantially different policies, whereas 9 have nearly identical policy texts}. Of the 116 apps, only 54\,\% had a GDPR-aware English policy \pc{arora2022}{7}{only 54\% of apps had English privacy policies that were “GDPR-aware”}.

\cf{Limits (card)} Data-practice scheme, not harms or stance; Google Play only; collected in July 2020. Agreement is among law students after training and weekly calibration, i.e.\ a favourable setting. Policy pairs were matched by app, not by provider.

\cf{Use for this thesis} (1) A realistic ceiling for agreement on policy text: trained law students reached only moderate $\kappa$, lower in German for most categories. Thresholds for D8 and D9 should be set against this, not against $\kappa$\,=\,0.8. (2) \textbf{Language rule (D7):} the German and English policies of the same app can differ, and the GDPR-specific text may exist only in the German one. ``English first, German fallback'' may therefore code the version written for non-EU users. Record both versions where they exist and test the difference in the pilot. (3) MAPP's German policies are a ready test set for prompts in German before the pilot. (4) The German-IP and browser-locale procedure is a precedent for the storefront settings of stage A.
\end{papercard}
